\scanpage{48}
% archive sheet number: 47

\deleted{\unsure{20|sa}} $20)$ \uline{Anneau à car.}

loc. sur $\mathcal E$, $A$ est de car $p>0$ ou de car. $0$, i.e. la famille

\[
\begin{aligned}
 &\deleted{\ensuremath{e_{\mathcal E}\longleftarrow\unsure{\mathop{\mathrm{Sup}}\limits_p|{\bigcup}_p}V(p\,1)}}\\
 &\text{des}\qquad
 \left\{\begin{array}{l}
 \deleted{\ensuremath{\unsure{V|p\in P}}}\ V(p.1_A)\hookrightarrow e_{\mathcal E}\\
 \displaystyle\bigcap_{n\geq 1}(e_{\mathcal E})_{n1_A}\lhook\joinrel\nearrow
 \end{array}\right.
\end{aligned}
\]

est épimorphique (cela fait \margin{\deleted{ou encore} $\mathcal E\longrightarrow\widehat R$ se factorise par le sous-topos $\mathrm{Sup}.$} intervenir une $\varprojlim$ \uline{infinie} \unsure{$\bigvee$|$\exists$}, donc ne donne pas l'existence d'un topos classifiant).

Ou encore $\mathcal E\to\widehat R$ se factorise localement sur $\mathcal E$ par la ``réunion'' $V=\deleted{\ensuremath{X}}_{\mathrm{car}>0}$ des sous-topos ouverts \deleted{\ensuremath{V}} $\widehat{R_{\mathbb Z/p\mathbb Z}}\simeq\widehat R_{\mathrm{Sp}\,\mathbb Z/p\mathbb Z}$, et le sous-topos $W=\widehat{R_{\mathbb Q}}$ (qui n'est pas un sous-topos ouvert) avec mêmes images \deleted{inverses}

\[
 X_{\mathrm{car}.}=\widetilde{R_{\mathrm{top.car}}}
\]

\deleted{\unsure{Or|Et}} \unsure{Or|On|Où} ``top car'' est la borne inférieure des ``top car $p$'' pour tt nb premier $p$, et pour $0$, i.e. la famille $X_i\longrightarrow X$ est couvrante ssi pour \inserted{$\forall p$} \deleted{\unsure{$p$, il|il}} \inserted{$y$ inclus $0$} existe un indice $i(p)$ tel que $(X_{i(p)})_p\longrightarrow (X)_p$ (fibre en $p$) ait une section. C'est

(p.ex. la famille des $(X_p\to X,\ p\ \text{premier ou nul})$ est couvrante, et si $X_{\mathbb Q}\ne\varnothing$, il n'y a pas des familles couvrantes plus \unsure{raffinées|raisonnables|cofinales})

\figure{\begin{tikzpicture}[x=1cm,y=1cm]
\node[draw,ellipse,inner sep=2pt] at (2.4,2.3) {Topos non \unsure{quasicompact|quasi-compact|cohérent}};
\draw (1.7,.85) ellipse[x radius=1.42,y radius=.76,rotate=-6];
\draw (3.03,1.02) ellipse[x radius=1.40,y radius=.82,rotate=23];
\draw[->] (2.0,1.34) to[bend right=15] (2.27,1.78);
\node at (2.13,1.93) {$W$}; \node at (1.00,1.07) {$Z'$};
\node at (3.34,1.14) {$Z''$}; \node at (2.13,.81) {$I$};
\draw (.85,.63) .. controls (1.2,.91) and (1.47,.55) .. (1.68,.66);
\draw (2.52,.67) .. controls (2.79,.92) and (2.99,.66) .. (3.20,.77);
\draw (.95,.16) to[bend right=15] (2.11,-.06) to[bend right=15] (3.20,.28);
\node at (2.05,-.33) {$V$};
\end{tikzpicture}}

Plus général\textsuperscript{t}, soit $B$ un topos, $V$ et $W$ deux sous-topos, considérons $T=T_B$ \deleted{\unsure{q|f}} et $T_V$ et $T_W$, enfin $T'=\{\mathcal E\to\{f\colon\mathcal E\to B$, qui se factorise \uline{local\textsuperscript{ment}} par $B$ \uline{ou} $B'$. Cela si\deleted{\unsure{gnifie|gnifie donc}} \unsure{aussi|ainsi}. Cela signifie que : \deleted{l'image} $e_{\mathcal E}=\mathrm{Sup}(e',e'')$, où $e'$ est le plus grand ouvert de $\mathcal E$ \deleted{sur lequel} \inserted{tel que} $f$ t. $\mathcal E_{/e'}$ se factorise par \unsure{$V$|$0$}, et idem pour $e''$. On a alors $f_{B/e'\wedge e''}\longrightarrow V\cap W$, \deleted{et on doit vérifier que} $e'\wedge e''$ est le plus gd ouvert \deleted{\ensuremath{\unsure{A|a}}} $\subset e_{\mathcal E}$ ayant cette propriété ? [si \deleted{si cela} \deleted{\ensuremath{\unsure{B|Z|U}}} est ains\textsuperscript{i}, on a $U\subset e'$, $U\subset e''$, d'où $U\subset e'\wedge e''$, OK.] /

% slanted
\margin{la factoris\textsuperscript{on} d'un \unsure{morphisme|élément} \unsure{par un ss-topos fermé|par ss-topos fermé} ou un ss-topos \unsure{ouvert|un} \unsure{d'une immersion|immersion} !}

S'il y a un topos classifiant pour $T'$, c'est un sous-topos \deleted{\ensuremath{\unsure{A|Z}}} $Z\subset B$ tel que $Z$ \deleted{admet un sous-topos} contenant $V$ et $W$, et ayant des sous-topos ouverts $Z'=Z_{/e'}$, $Z''=Z_{/e''}$, avec $Z'\subset V$, $Z''\subset W$, et \deleted{\ensuremath{\unsure{\mathcal E|Z}}} $\mathrm{Sup}(e',e'')=e_Z$ i.e. $Z=\mathrm{Sup}(Z',Z'')$. (dans l'ens des sous-topos de $Z$, ou de $B$, c'est kif kif). On \unsure{aurait|aurai|aura} donc

\[
 Z=\mathrm{Sup}(Z',Z'')\subset\mathrm{Sup}(V,W)\subset Z
\]

donc $Z=\mathrm{Sup}(V,W)$. De plus on aura \deleted{(ou $Z'$ sera \unsure{l'intersection|l'image} de $Z'$.}

\[
 \begin{aligned}
 \deleted{\text{\unsure{nécessairement|nécessaires}}}\quad
 \deleted{\ensuremath{\unsure{V\times Z|V\times_Z}}}\quad
 Z'&=Z'\cap V'=Z'\times_Z V=\text{ouvert de }V\\
 \unsure{\text{ou }=|\text{ou }-}\qquad Z''&=\text{ouvert de }W
 \end{aligned}
\]
